% =====================================================
% Chapter 4: System Architecture
% =====================================================
\chapter{System Architecture}
\label{ch:architecture}

This chapter presents the system-level architecture of an AGC microphone preamplifier in general/conceptual terms, before the specific circuit realisation is discussed in Chapter~\ref{ch:circuitdesign}. The block diagrams in this chapter illustrate the functional signal flow common to this class of circuit and are \textbf{conceptual reference diagrams only}: they do not represent the team's final implemented circuit, and no component values or implementation-specific details are assigned to them. The team's actual implemented block diagram, once finalised, should replace or supplement these figures.

\section{Overall System Block Diagram}

Figure~\ref{fig:system-block} shows the overall signal path from the microphone to the system output, including the AGC feedback loop as a single functional block. Each block represents a function that must be realised by the final circuit; the specific stage-by-stage implementation is presented in Chapter~\ref{ch:circuitdesign}.

% =====================================================
% CONCEPTUAL / REFERENCE DIAGRAM ONLY.
% This is a generic AGC microphone preamplifier system block
% diagram used for explanatory purposes in Chapter 4. It does
% NOT represent the team's actual implemented circuit and
% carries no component values or implementation-specific detail.
% Replace/supplement with the team's real system diagram once
% the design is finalised.
% =====================================================
\begin{figure}[H]
\centering
\begin{tikzpicture}[
    node distance=1.1cm and 1.3cm,
    block/.style={rectangle, draw, thick, minimum height=1.1cm, minimum width=2.6cm, align=center, fill=gray!8},
    every node/.style={font=\small}
]
    \node[block] (mic) {Microphone\\(input transducer)};
    \node[block, right=of mic] (pre) {Input buffer /\\preamplifier stage};
    \node[block, right=of pre, minimum width=3.4cm] (vga) {Variable-gain stage\\with AGC feedback};
    \node[block, right=of vga] (out) {Output buffer /\\line-level output};

    \draw[-{Latex}, thick] (mic) -- (pre);
    \draw[-{Latex}, thick] (pre) -- (vga);
    \draw[-{Latex}, thick] (vga) -- (out);

    \node[below=0.7cm of vga, font=\footnotesize\itshape] (note) {(AGC feedback loop detailed in Fig.~\ref{fig:agc-block})};
    \draw[-{Latex}, thick, dashed] (vga.south) -- (note.north);
\end{tikzpicture}
\caption{Conceptual (reference) system-level block diagram of an AGC microphone preamplifier. \textbf{This diagram is generic and illustrative; it is not the project's final implemented circuit.}}
\label{fig:system-block}
\end{figure}


\section{AGC Feedback Loop -- Detailed Block Diagram}

Figure~\ref{fig:agc-block} expands the ``Variable-Gain Stage with AGC Feedback'' block of Figure~\ref{fig:system-block} into its constituent functional blocks, following the general three-function AGC structure introduced in Chapter~\ref{ch:theory} (level detection, control-voltage generation, and gain adjustment).

% =====================================================
% CONCEPTUAL / REFERENCE DIAGRAM ONLY.
% This is a generic AGC feedback-loop block diagram used for
% explanatory purposes in Chapter 4. It illustrates the
% textbook-level structure of an AGC loop (level detection,
% control-voltage generation, gain adjustment) and does NOT
% represent the team's actual implemented circuit. No component
% values or implementation-specific detail are assigned here.
% =====================================================
\begin{figure}[H]
\centering
\begin{tikzpicture}[
    node distance=1.0cm and 1.2cm,
    block/.style={rectangle, draw, thick, minimum height=1.0cm, minimum width=2.5cm, align=center, fill=gray!8},
    sum/.style={circle, draw, thick, minimum size=0.7cm},
    every node/.style={font=\small}
]
    \node[block] (vga) {Variable-Gain\\Amplifier (VGA)};
    \node[block, right=1.6cm of vga] (det) {Level / Envelope\\Detector};
    \node[sum, below=1.6cm of det] (cmp) {$\Sigma$};
    \node[block, below=1.6cm of vga, minimum width=3.0cm] (ctrl) {Control-Voltage\\Generator (Attack/Release)};

    \node[left=1.2cm of vga, font=\footnotesize] (in) {Input};
    \node[right=1.2cm of det, font=\footnotesize] (out) {Output};

    \draw[-{Latex}, thick] (in) -- (vga);
    \draw[-{Latex}, thick] (vga) -- (det);
    \draw[-{Latex}, thick] (det) -- (out);
    \draw[-{Latex}, thick] (det.south) -- (cmp.north);
    \node[below=0.55cm of cmp, font=\footnotesize, align=center] (ref) {Reference\\threshold};
    \draw[-{Latex}, thick] (ref) -- (cmp);
    \draw[-{Latex}, thick] (cmp) -- (ctrl.east);
    \draw[-{Latex}, thick] (ctrl) -- (vga.south) node[midway, left, font=\footnotesize\itshape] {control voltage};
\end{tikzpicture}
\caption{Conceptual (reference) AGC feedback-loop block diagram, following the general three-function structure (level detection, control-voltage generation, gain adjustment) described in Chapter~\ref{ch:theory}. \textbf{This diagram is generic and illustrative; it is not the project's final implemented circuit.}}
\label{fig:agc-block}
\end{figure}


\section{Description of Functional Blocks}

\begin{itemize}
    \item \textbf{Microphone (input transducer)} -- converts acoustic pressure to a small-signal AC voltage; treated as a Th\'{e}venin source with its own output impedance and noise characteristics.
    \item \textbf{Input Buffer / Preamplifier Stage} -- provides high input impedance and initial low-noise gain, setting up the signal for the variable-gain stage.
    \item \textbf{Variable-Gain Stage} -- an amplifier stage whose gain can be varied electronically by a control signal (see Chapter~\ref{ch:theory} for candidate implementations).
    \item \textbf{Level/Envelope Detector} -- estimates the magnitude of the (typically output) signal over a short time window.
    \item \textbf{Reference/Threshold Comparison and Control Voltage Generator} -- compares the detected level with a target threshold and produces a smoothed DC control signal, shaped by the attack/release filtering discussed in Chapter~\ref{ch:theory}.
    \item \textbf{Output Buffer} -- isolates the variable-gain stage from the load and delivers the final output signal.
\end{itemize}

\section{Mapping to the Implemented Circuit}

\placeholder{Once the team's final schematic is available, add a short paragraph here mapping each block above to the corresponding stage(s) in the schematic of Chapter~\ref{ch:circuitdesign} (e.g. ``the Variable-Gain Stage is realised using an op-amp with a JFET in its feedback path, Q1'')}
